
\section{Spherical tensors and Wigner--Eckart}
A rank-\(k\) irreducible spherical tensor \(T_q^{(k)}\) has
\(2k+1\) components \(q=-k,\dots,k\). Rank 1 behaves as a vector;
rank 2 has five quadrupole-like components. The Wigner--Eckart theorem
separates angular geometry from dynamics:
\[
\boxed{\langle j'm'|T_q^{(k)}|jm\rangle
=(-1)^{j'-m'}
\begin{pmatrix}j'&k&j\\-m'&q&m\end{pmatrix}
\langle j'\Vert T^{(k)}\Vert j\rangle.}
\]
The \(3j\) symbol enforces \(m'=m+q\), the triangle rule
\(|j-j'|\le k\le j+j'\), and allowed integer/half-integer sums.
A nonzero \(3j\) does \emph{not} guarantee that the reduced matrix
element, and hence the transition, is nonzero.

\subsection{A simple \(3j\) symbol}
Adding angular momentum zero changes nothing:
\(\langle0,0;1,0|1,0\rangle=1\). Hence
\[
\begin{pmatrix}0&1&1\\0&0&0\end{pmatrix}
=\frac{(-1)^{0-1+0}}{\sqrt{2(1)+1}}
\langle0,0;1,0|1,0\rangle=-\frac1{\sqrt3}.
\]

\section{Electromagnetic multipoles and parity}
Electric multipole \(Ek\) has rank \(k\), parity \((-1)^k\).
Magnetic multipole \(Mk\) has rank \(k\), parity \((-1)^{k+1}\).
Both require the angular-momentum triangle rule and projection rule.

\begin{center}
\begin{tabularx}{\linewidth}{@{}lclX@{}}
\toprule
Type & Rank & Operator parity & Common orbital screening (single electron)\\
\midrule
E1 & 1 & odd & Opposite-parity orbitals; \(\Delta l=\pm1\)\\
M1 & 1 & even & Same parity; commonly \(\Delta l=0\)\\
E2 & 2 & even & Same parity; \(\Delta l=0,\pm2\), plus triangle rule\\
\bottomrule
\end{tabularx}
\end{center}
For E1, \(\Delta J=0,\pm1\) but \(J=0\leftrightarrow J'=0\)
is forbidden. For E2, \(|J'-J|\le2\le J'+J\); listing
\(\Delta J=0,\pm1,\pm2\) without this triangle condition is incomplete.
\begin{methodbox}
\textbf{Exam method:} identify \(Ek/Mk\); check parity, triangle and
projection rules; apply the specified spin/orbital assumptions;
then calculate the \(3j\) and reduced/radial elements.
\end{methodbox}
